% ============================================================================
% METHODS BRIEF: POLYPHARMACY RISK ASSESSMENT
% Publication Format
% ============================================================================

\documentclass[11pt,a4paper]{article}
\usepackage[utf8]{inputenc}
\usepackage{amsmath,amssymb}
\usepackage{booktabs}
\usepackage{geometry}
\usepackage{graphicx}
\usepackage{hyperref}
\usepackage{float}
\geometry{margin=1in}

\begin{document}

\section*{Polypharmacy Risk Assessment Methods}

\subsection*{Polypharmacy Risk Index (PRI)}

We developed the Polypharmacy Risk Index (PRI) to quantify individual drug risk within multi-drug regimens. The PRI combines four network-based metrics weighted according to their clinical relevance:

\begin{equation}
\text{PRI}(d) = 0.25 \cdot C_{\text{degree}}(d) + 0.30 \cdot C_{\text{weighted}}(d) + 0.20 \cdot C_{\text{betweenness}}(d) + 0.25 \cdot S(d)
\end{equation}

\subsection*{Component Metrics}

\begin{itemize}
    \item \textbf{Degree Centrality} $C_{\text{degree}}(d)$: Captures the interaction burden as the number of drugs interacting with drug $d$.
    \item \textbf{Weighted Degree} $C_{\text{weighted}}(d)$: Represents the severity-weighted sum of all interactions, giving greater importance to contraindicated and major interactions.
    \item \textbf{Betweenness Centrality} $C_{\text{betweenness}}(d)$: Quantifies the drug's role in risk propagation pathways across the network.
    \item \textbf{Severity Profile Score} $S(d)$: Measures the proportion of a drug's interactions that are classified as contraindicated or major.
\end{itemize}

\subsection*{Risk Classification Thresholds}

\begin{table}[H]
\centering
\caption{PRI Risk Classification Thresholds}
\label{tab:pri_thresholds}
\begin{tabular}{lll}
\toprule
\textbf{PRI Score} & \textbf{Risk Level} & \textbf{Clinical Action} \\
\midrule
$> 0.5$ & High Risk & Immediate clinical review required \\
$0.3 - 0.5$ & Medium Risk & Close monitoring warranted \\
$< 0.3$ & Lower Risk & Standard monitoring protocols \\
\bottomrule
\end{tabular}
\end{table}

\subsection*{Severity Weighting Schema}

The severity weights used in the weighted degree calculation are derived from clinical significance:

\begin{table}[H]
\centering
\caption{Severity Weights for PRI Calculation}
\label{tab:severity_weights}
\begin{tabular}{lr}
\toprule
\textbf{Severity Level} & \textbf{Weight} \\
\midrule
Contraindicated & 4.0 \\
Major & 3.0 \\
Moderate & 2.0 \\
Minor & 1.0 \\
\bottomrule
\end{tabular}
\end{table}

\subsection*{DDI Severity Distribution}

The DDI severity distribution after recalibration using the empirical keyword-based classifier validated against DDInter (66.4\% exact accuracy, $\kappa$ = +0.096):

\begin{table}[H]
\centering
\caption{DDI Severity Distribution}
\label{tab:severity_distribution}
\begin{tabular}{lrr}
\toprule
\textbf{Severity} & \textbf{Count} & \textbf{Percentage} \\
\midrule
Moderate & 608,742 & 80.1\% \\
Minor & 70,682 & 9.3\% \\
Contraindicated & 44,306 & 5.8\% \\
Major & 36,044 & 4.7\% \\
\midrule
\textbf{Total} & \textbf{759,774} & \textbf{100.0\%} \\
\bottomrule
\end{tabular}
\end{table}

\section*{Figures}

\begin{figure}[H]
\centering
\includegraphics[width=0.8\textwidth]{plots_and_data/polypharmacy_risk_escalation.png}
\caption{Polypharmacy risk escalation visualization showing how overall regimen risk increases with the number of concurrent medications.}
\label{fig:risk_escalation}
\end{figure}

\begin{figure}[H]
\centering
\includegraphics[width=0.8\textwidth]{plots_and_data/drug_risk_matrix.png}
\caption{Drug risk matrix displaying high-risk drug combinations and their associated severity levels.}
\label{fig:risk_matrix}
\end{figure}

\begin{figure}[H]
\centering
\includegraphics[width=0.8\textwidth]{plots_and_data/severity_heatmap.png}
\caption{Severity heatmap showing the distribution of interaction severities across drug classes.}
\label{fig:severity_heatmap}
\end{figure}

\bibliographystyle{plain}
\begin{thebibliography}{2}
\bibitem{wishart2018drugbank} Wishart DS, et al. DrugBank 5.0: a major update to the DrugBank database for 2018. \textit{Nucleic Acids Res}. 2018;46(D1):D1074-D1082.
\bibitem{xiong2022ddinter} Xiong G, et al. DDInter: an online drug-drug interaction database. \textit{Nucleic Acids Res}. 2022;50(D1):D1200-D1207.
\end{thebibliography}

\end{document}
